\section{Discussion}

This study supports a data-centric view of cancer multi-omics prediction. In the five TCGA cancer-internal tasks, the same nominal multi-omics feature table behaved differently depending on modality coverage and sample pairing. The sample-level audit confirmed that the starting tables were uniquely keyed by labelled tumor sample, with no duplicate sample identifiers and high core-modality completeness. This is important because the subsequent mismatch results can be interpreted as a stress test of broken sample pairing rather than a correction for obvious duplicate rows. RPPA remained sparse at the feature level and should not be treated as a routine early-fusion modality in these tasks. More importantly, simulated sample mismatch degraded performance even though each modality's marginal distribution remained intact. This is exactly the kind of failure that can be invisible to feature-level summary statistics and visible only after sample-level alignment is audited.

The mismatch stress test also shows why a single model score can be misleading. UCEC and COADREAD were strongly affected by non-mRNA modality mismatch, especially for ExtraTrees and the Liquid-family models, whereas KIRC and PRAD were less sensitive. This does not mean that alignment is unimportant for KIRC or PRAD; rather, it means that these particular endpoints and feature sets were less dependent on cross-modal relationships under the tested models. A useful multi-omics benchmark should therefore report both performance and alignment sensitivity for classical and neural model families.

The BRCA external-validation case study adds a complementary perspective. CPTAC 2020 was more separable from the TCGA/METABRIC source distribution than SMC 2018 or SCAN-B/GSE96058 in marker expression space, yet external macro F1 depended on the model and cohort. This finding cautions against interpreting domain shift as a single scalar quantity. Domain classifier separability, standardized mean shift, class balance, label definitions, and assay compatibility should be considered together.

The 10-cancer exploratory extension clarifies a common pitfall in multi-cancer machine learning. Pan-cancer cancer-type prediction can look nearly solved, but that success mostly reflects strong tissue or cancer identity signals. It does not imply that clinically meaningful within-cancer endpoints are equally easy. The endpoint screen identified 21 candidate within-cancer tasks, but their apparent difficulty varied sharply by endpoint family: subtype tasks were generally learnable, grade tasks were cancer-dependent, and stage or T-stage tasks were often weak molecular prediction targets. Therefore, multi-cancer benchmarks should include within-cancer endpoints and should not rely only on cancer-type classification as evidence of biomedical predictive utility.

The expanded 10-fold baseline benchmark further argues against overinterpreting small model-rank differences. Although tree ensembles, logistic regression, and neural architectures each performed well in different tasks, no model family was consistently dominant across all endpoints. This is particularly relevant for Liquid/CfC-style models: their inclusion is useful as a modern neural comparator, but the present data do not justify an architecture-superiority claim. Instead, the results suggest an applicability boundary. Liquid/CfC was competitive for COADREAD, UCEC, and PRAD, but was weaker for HNSC HPV status and KIRC grade. This pattern is biologically and statistically plausible: when a simpler, high-signal decision boundary exists, regularized linear models or tree ensembles can be sufficient; when cross-modality relationships are more complex, Liquid/CfC-style modality-sequence models may be competitive. The medical interpretation table strengthens this view by linking model behavior to endpoint biology rather than treating all tasks as interchangeable labels. The stronger conclusion is methodological: benchmark studies should make sample alignment, modality coverage, endpoint biology, cohort shift, and robustness checks part of the evaluation protocol.

The main limitation is that the mismatch experiment is simulated. Real sample misalignment can arise through more complex mechanisms, including patient/sample ID mismatches, aliquot-level differences, assay-specific filtering, clinical label harmonization errors, and differences between sample, portion, analyte, and aliquot provenance. Our audit was conducted at the harmonized cBioPortal sample-identifier level; it verifies one labelled modelling row per tumor sample in the analysed tables, but it is not equivalent to a complete aliquot-level molecular provenance audit. The simulation is therefore best interpreted as a stress test of sensitivity, not as an estimate of the frequency of real errors. A second limitation is that external validation was performed as a BRCA case study. Extending the domain-shift component to external cohorts for UCEC, COADREAD, HNSC, KIRC, and PRAD would further strengthen the benchmark. A third limitation is that the 21-task endpoint screen used lightweight baselines and three-fold cross-validation for scope discovery; because endpoint discovery and apparent difficulty estimation used the same public resource, these results are subject to selection bias. Future work should apply the full 10-fold baseline, mismatch, and external-validation protocol to the most clinically important tasks identified by this screen.

Despite these limitations, the second manuscript is distinct from a model-comparison paper. Its central claim is not that one architecture is best, but that data reliability checks should be part of the benchmark itself. We recommend that cancer multi-omics prediction studies report: the sample alignment key, duplicate sample checks, modality coverage by labelled sample, coverage differences by cancer type, compatible external-cohort modality sets, training-only preprocessing, one-vs-rest domain or cancer-identity separability, and mismatch or permutation stress tests.
